\section{Технологический раздел}
В данном разделе рассматриваются различные реляционные и 
графовые СУБД, проводится их сравнение по определенным критериям.
Также описываются выбранные средства реализации и тестирование, 
демонстрируются результаты тестирования и функциональность разработанного
приложения.

\subsection{Выбор СУБД}
В соответствии с аналитическим разделом для хранения данных разрабатываемой 
системы необходимо использовать две различные систем управления базами 
данных: реляционную и графовую.

\subsubsection{Выбор реляционной СУБД}
Для хранения информации о пользователях, ролях, командах и конфигурациях 
систем необходима реляционная система управления базами данных. Наиболее 
распространенными решениями в данной области являются:
\begin{itemize}
    \item PostgreSQL~\cite{};
    \item MySQL~\cite{};
    \item Microsoft SQL Server~\cite{};
    \item Oracle Database~\cite{}.
\end{itemize}

В таблице \ref{tbl:relational-dbms} приводится сравнение рассматриваемых 
реляционных систем по выделенным критериям.

\begin{center}
    \begin{tabularx}{\textwidth}{|X|X|X|X|}
        \caption{\label{tbl:relational-dbms}Сравнение реляционных систем управления базами данных} \\ \hline
        Решение & 
        Открытый исходный код и бесплатное распространение & 
        Поддержка стандарта SQL & 
        Кроссплатфор-менность и поддержка запуска в изолированных контейнерах \\ \hline
        \endfirsthead

        \multicolumn{4}{l}{{Продолжение таблицы\ \thetable{}}} \\ \hline
        Решение & 
        Открытый исходный код и бесплатное распространение & 
        Поддержка стандарта SQL & 
        Кроссплатфор-менность и поддержка запуска в изолированных контейнерах \\ \hline
        \endhead

        \endfoot
        \endlastfoot

        PostgreSQL & Да & Да & Да \\ \hline
        MySQL & Да & Частично & Да \\ \hline
        Microsoft SQL Server & Нет & Да & Да \\ \hline
        Oracle Database & Нет & Да & Да \\ \hline
    \end{tabularx}
\end{center}

В качестве используемой реляционной системы управления базами данных был выбран 
PostgreSQL, так как он имеет полностью открытый код, распространяется бесплатно
и полностью поддерживает стандарт SQL.

\subsubsection{Выбор графовой СУБД}
Для хранения топологии инфраструктуры и выполнения алгоритмического обхода 
графа необходима графовая система управления базами данных. Наиболее 
популярными решениями на рынке являются:
\begin{itemize}
    \item Neo4j~\cite{};
    \item ArangoDB~\cite{};
    \item OrientDB~\cite{}.
\end{itemize}

В таблице \ref{tbl:graph-dbms} приводится сравнение рассматриваемых графовых систем по выделенным критериям.

\begin{center}
    \begin{tabularx}{\textwidth}{|X|X|X|X|}
        \caption{\label{tbl:graph-dbms}Сравнение графовых систем управления базами данных} \\ \hline
        Решение & 
        Нативная графовая архитектура хранения и обработки данных & 
        Наличие встроенной библиотеки графовых алгоритмов & 
        Наличие бесплатной версии \\ \hline
        \endfirsthead

        \multicolumn{4}{l}{{Продолжение таблицы\ \thetable{}}} \\ \hline
        Решение & 
        Нативная графовая архитектура хранения и обработки данных & 
        Наличие встроенной библиотеки графовых алгоритмов & 
        Наличие бесплатной версии \\ \hline
        \endhead

        \endfoot
        \endlastfoot

        Neo4j & Да & Да & Да \\ \hline
        ArangoDB & Нет & Нет & Да \\ \hline
        OrientDB & Нет & Нет & Да \\ \hline
    \end{tabularx}
\end{center}

Решения ArangoDB и OrientDB являются мультимодельными базами данных, 
совмещающими в себе документо-ориентированный и графовый подходы, что 
приводит к снижению производительности при глубоком рекурсивном обходе связей. 
СУБД Neo4j является полностью нативной графовой базой данных и предоставляет 
доступ к специализированной библиотеке графовых алгоритмов, идеально 
подходящей для вычисления единых точек отказа и выявления циклических 
зависимостей. По результатам сравнения в качестве графовой системы управления 
базами данных была выбрана Neo4j.

\subsection{Выбор средств реализации приложения}
В качестве языка программирования для разработки серверной части 
приложения был выбран C\#~\cite{}. Данный язык обладает строгой статической типизацией, 
богатым набором встроенных структур данных и поддержкой парадигм 
объектно-ориентированного и асинхронного программирования. Разработка ведется 
на базе современной кроссплатформенной платформы .NET 10~\cite{}. Выбор данной 
платформы обусловлен ее высокой производительностью, встроенной поддержкой 
внедрения зависимостей и развитой экосистемой, что делает ее оптимальным 
решением для создания надежных и масштабируемых серверных приложений.

Для взаимодействия с реляционной базой данных PostgreSQL была выбрана технология 
объектно-реляционного отображения Entity Framework Core~\cite{}. Использование 
EF Core позволяет абстрагироваться от написания низкоуровневых SQL-запросов 
и обеспечивает манипулирование данными посредством строго типизированных LINQ-запросов, 
а также автоматизируя процесс управления схемой базы данных через механизм миграций. 
В свою очередь, для интеграции с графовой базой данных Neo4j применяется официальный 
драйвер Neo4j.Driver~\cite{}. Такое разделение средств доступа к данным позволяет 
эффективно реализовать паттерн репозитория, инкапсулируя специфику работы с 
каждой конкретной СУБД.

\subsection{Создание объектов БД}
Для обеих PostgreSQL процесс создания таблиц и настройки ограничений возложен
на используемый фреймворк EF Core PostgreSQL. В графовой СУБД Neo4j таблиц нет,
а структура всегда одинакова: узлы и связи между ними, поэтому описание структуры
заранее не производится.

\subsubsection{Создание объектов реляционной БД}
Фреймворк EF Core использует механизм миграций. Миграция БД --- это процесс 
переноса данных, структур или кода из одной среды, системы или формата в другую. 
В разработке механизм миграций используется как <<версионный контроль>> структуры 
БД с помощью кода.

Структура БД в фреймворке EF Core описывается на выбранном языке программирования.
Согласно описанию фреймворк формирует SQL запросы для создания описанных объектов, 
вспомогательных таблиц для отслеживания миграций, ограничений, индексов.
Каждое изменение структуры БД оформляется в виде отдельной миграции. Исходный код
описания структуры разрабатываемой БД представлен в листинге~\ref{lst:relational-init}.

\begin{listing}[Описание структуры реляционной БД]{lst:relational-init}{csh}
public class AnalyzerDbContext(DbContextOptions<AnalyzerDbContext> options) : DbContext(options)
{
    public DbSet<User> Users => Set<User>();
    public DbSet<Team> Teams => Set<Team>();
    public DbSet<Invite> Invites => Set<Invite>();
    public DbSet<ITSystem> ITSystems => Set<ITSystem>();

    protected override void OnModelCreating(ModelBuilder modelBuilder)
    {
        base.OnModelCreating(modelBuilder);

        modelBuilder.Entity<User>(builder =>
        {
            builder.HasKey(u => u.Id);
            builder.HasIndex(u => u.Username).IsUnique();
            builder.HasIndex(u => u.Email).IsUnique();

            builder.Property(u => u.Role).HasConversion<int>();

            builder.HasOne<Team>()
                   .WithMany()
                   .HasForeignKey(u => u.TeamId)
                   .OnDelete(DeleteBehavior.Restrict);
        });

        modelBuilder.Entity<Team>(builder =>
        {
            builder.HasKey(t => t.Id);

            builder.Ignore(t => t.MemberIds);
            builder.Ignore("_memberIds");
        });

        modelBuilder.Entity<Invite>(builder =>
        {
            builder.HasKey(i => i.Id);
            builder.HasIndex(i => i.Code).IsUnique();

            builder.Property(i => i.Status).HasConversion<int>();
            builder.Property(i => i.Role).HasConversion<int>();

            builder.HasOne<Team>()
                   .WithMany()
                   .HasForeignKey(i => i.TeamId)
                   .OnDelete(DeleteBehavior.Cascade);

            builder.HasOne<User>()
                   .WithMany()
                   .HasForeignKey(i => i.ActivatedByUserId)
                   .OnDelete(DeleteBehavior.SetNull);
        });

        modelBuilder.Entity<ITSystem>(builder =>
        {
            builder.HasKey(s => s.Id);
            builder.HasIndex(s => new { s.TeamId, s.Name }).IsUnique();

            builder.HasOne<Team>()
                   .WithMany()
                   .HasForeignKey(s => s.TeamId)
                   .OnDelete(DeleteBehavior.Cascade);
        });
    }
}
\end{listing}

\subsection{Разработка запросов к графовой БД}